% Original analytic teaching diagram, 2026-11-08.
\definecolor{VB}{HTML}{7188A3}\definecolor{VBF}{HTML}{E5EAF0}
\definecolor{VG}{HTML}{829C8B}\definecolor{VGF}{HTML}{E8EEE9}
\definecolor{VY}{HTML}{B4A16F}\definecolor{VYF}{HTML}{F1ECDD}
\definecolor{VR}{HTML}{A7787D}\definecolor{VRF}{HTML}{F0E5E6}
\definecolor{VP}{HTML}{9588A5}\definecolor{VPF}{HTML}{ECE8F0}
\definecolor{VI}{HTML}{3B4252}\definecolor{VM}{HTML}{6B7280}\definecolor{VL}{HTML}{4C4D4F}
% Phi(x)=(x,x)，V=span{(1,1)}；参数 (2,0),(1,1),(0,2) 都给出 f(x)=2x（虚线 w1+w2=2）。比例 1 单位=17mm。
\begin{tikzpicture}[x=1mm,y=1mm,>=Stealth,font=\figurefont\fontsize{8}{10}\selectfont,
lab/.style={text=VM,font=\figurefont\fontsize{7.5}{9}\selectfont},
ax/.style={->,VL,line width=.65pt},guide/.style={VM,line width=.55pt,dash pattern=on 1.6pt off 1.2pt}]
\def\u{17}
\draw[ax] (-.45*\u,0)--(2.65*\u,0) node[right=-.2mm,text=VI] {$w_1$};
\draw[ax] (0,-.45*\u)--(0,2.6*\u) node[above=-.2mm,text=VI] {$w_2$};
\draw[VB,line width=1.1pt] (-.4*\u,-.4*\u)--(2.15*\u,2.15*\u) node[above right=-.3mm,text=VI] {$V$};
\draw[VY,line width=1pt,dash pattern=on 2.2pt off 1.4pt] (-.28*\u,2.28*\u)--(2.28*\u,-.28*\u);
\draw[VL,line width=.5pt] (.9*\u,.9*\u)--(.97*\u,.83*\u)--(1.07*\u,.93*\u);
\draw[VR,->,line width=.85pt] (0,0)--(2*\u,0);
\draw[VR,->,line width=.85pt] (0,0)--(0,2*\u);
\draw[VG,->,line width=1.2pt] (0,0)--(1*\u,1*\u);
\fill[VR] (2*\u,0) circle (.9);\fill[VR] (0,2*\u) circle (.9);\fill[VG] (\u,\u) circle (1);
\node[below=.6mm,text=VI] at (2*\u,0) {$(2,0)$};
\node[right=.5mm,text=VI] at (0,2*\u) {$(0,2)$};
\node[right=1.6mm,text=VI] at (\u,\u) {$(1,1)$};
\end{tikzpicture}
